\chapter{T-Duality and the Self-Dual Radius}\label{ch:tduality}

In \cref{ch:circle} we compactified one direction of space on a circle of radius $R$ and found
that a closed string carries two integers there: a momentum number $n$, as any quantum particle
on a circle does, and a winding number $w$, which only an extended object can have. This
chapter asks what happens when the circle is made small. For a point particle the answer is
dull: the Kaluza--Klein momentum states become infinitely heavy, the circle disappears from
low-energy physics, and that is all. For a string the answer is that \emph{nothing singular
happens at all}: the theory on a circle of radius $R$ is the same theory as on a circle of
radius $\ap/R$, with the roles of momentum and winding exchanged. This statement is called
\emph{T-duality}\index{T-duality}. It is the simplest example of a phenomenon with no
counterpart in field theory, it is the seed from which the duality group $\Odd{d}$, the
doubled coordinates and the generalized metric of Parts~II and~III grow, and it is the origin
of the ``two scales'' in the title of this book.

We shall be careful about one thing that is often blurred. ``T-duality is a symmetry'' can mean
four different statements of increasing strength: that the mass formula is invariant; that the
two conformal field theories are the same; that string perturbation theory agrees at every
loop order; and that the full theory is the same. The evidence, and the part that a machine
has checked, differ from level to level, and a student should know at each moment which level is
being discussed. \Cref{sec:tduality:levels} lays out this ladder; the rest of the chapter climbs it.

\section{The symmetry of the mass formula}\label{sec:tduality:mass}

\subsection{Data carried over from the circle}

We use the conventions of \cref{ch:circle}: $\ap$ has dimension length$^2$, the string
length is $\ell_s=\sqrt{\ap}$, the compact coordinate is $X\equiv X^{25}\sim X+2\pi R$, the
momentum along the circle is $n/R$ with $n\in\Z$, and the string may close only up to a winding,
$X(\sigma+2\pi)=X(\sigma)+2\pi wR$ with $w\in\Z$\index{winding number}\index{momentum number}. (Tong calls the winding number $m$; we write $w$
throughout.) The zero modes are packaged into the left- and right-moving momenta
\begin{equation}\label{eq:tduality:pLpR}
  p_L=\frac nR+\frac{wR}{\ap},\qquad p_R=\frac nR-\frac{wR}{\ap},
\end{equation}
which appear in the mode expansion $X=X_L(\sigma^+)+X_R(\sigma^-)$, $\sigma^\pm=\tau\pm\sigma$,
\begin{equation}\label{eq:tduality:modes}
\begin{aligned}
  X_L(\sigma^+)&=\tfrac12x+\tfrac12\ap\,p_L\,\sigma^+
     +\ii\sqrt{\tfrac{\ap}{2}}\sum_{k\ne0}\frac1k\,\tilde\alpha_k\,\ee^{-\ii k\sigma^+},\\
  X_R(\sigma^-)&=\tfrac12x+\tfrac12\ap\,p_R\,\sigma^-
     +\ii\sqrt{\tfrac{\ap}{2}}\sum_{k\ne0}\frac1k\,\alpha_k\,\ee^{-\ii k\sigma^-}
\end{aligned}
\end{equation}
\source{Tong §8.2, eq.~(8.4), ll.~11351--11415}. As a check, adding the two lines gives
$X=x+\frac12\ap(p_L+p_R)\tau+\frac12\ap(p_L-p_R)\sigma+\dots=x+\ap\frac nR\tau+wR\,\sigma+\dots$: the
centre of mass moves with momentum $n/R$ and the string winds $w$ times. The $L_0$ and
$\tilde L_0$ constraints give the mass seen by an observer in the 25 non-compact dimensions,
\begin{equation}\label{eq:tduality:M2LR}
  M^2=p_L^2+\frac4{\ap}(\tilde N-1)=p_R^2+\frac4{\ap}(N-1),
\end{equation}
where $N$ and $\tilde N$ are the oscillator levels of the $\alpha_k$ and the $\tilde\alpha_k$
\source{Tong §8.2, ll.~11430--11452}. The \emph{difference} of the two expressions is the
level-matching condition\index{level matching}, and their \emph{average} is the mass formula:
\begin{align}
  p_L^2-p_R^2=\frac{4nw}{\ap}\quad&\Longrightarrow\quad N-\tilde N=nw,
     \label{eq:tduality:level}\\
  M^2=\tfrac12(p_L^2+p_R^2)+\frac2{\ap}(N+\tilde N-2)
     &=\frac{n^2}{R^2}+\frac{w^2R^2}{\ap^2}+\frac2{\ap}(N+\tilde N-2)
     \label{eq:tduality:M2}
\end{align}
\source{Tong §8.2, eqs.~(8.5)--(8.6), ll.~11458--11480}. The two new terms have transparent
meanings. A state with momentum $n/R$ along a hidden direction looks, in the lower-dimensional
world, like a particle of mass $|n|/R$. A string wound $w$ times around a circle of circumference
$2\pi R$ has length $2\pi|w|R$ and therefore rest energy $2\pi|w|R\,T=|w|R/\ap$, with
$T=1/(2\pi\ap)$ the tension \source{Tong §8.2, ll.~11482--11487}.

\subsection{The duality map}

\begin{definition}\label{def:tduality:map}
The \emph{T-duality map} on the labels of closed-string states on a circle is
\begin{equation}\label{eq:tduality:map}
  \mathsf T:\quad R\longmapsto\tilde R=\frac{\ap}{R},\qquad (n,w)\longmapsto(\tilde n,\tilde w)=(w,n),
  \qquad (N,\tilde N)\longmapsto(N,\tilde N).
\end{equation}
\end{definition}

\begin{proposition}[T-duality of the zero modes]\label{thm:tduality:pLpR}
Under $\mathsf T$,
\begin{equation}\label{eq:tduality:pflip}
  p_L\longmapsto p_L,\qquad p_R\longmapsto-p_R .
\end{equation}
Consequently the mass formula \eqref{eq:tduality:M2} and the level-matching condition
\eqref{eq:tduality:level} are invariant: the states at radius $R$ and the states at radius
$\ap/R$ are in one-to-one correspondence, with equal masses.
\end{proposition}

\begin{proof}
Substitute \eqref{eq:tduality:map} into \eqref{eq:tduality:pLpR}:
\[
  \tilde p_L=\frac{\tilde n}{\tilde R}+\frac{\tilde w\tilde R}{\ap}
    =\frac{wR}{\ap}+\frac{n}{\ap}\cdot\frac{\ap}{R}=\frac{wR}{\ap}+\frac nR=p_L,\qquad
  \tilde p_R=\frac{wR}{\ap}-\frac nR=-p_R .
\]
The mass \eqref{eq:tduality:M2LR} contains only $p_L^2$, $p_R^2$, $N$ and $\tilde N$, so it is
unchanged. Level matching involves the product $nw=wn$, also unchanged. The map is a bijection
of $\Z^2$ because it is its own inverse.
\end{proof}

The rule \eqref{eq:tduality:pflip} is what Tong and Giveon--Porrati--Rabinovici state
\source{Tong §8.3, ll.~11627--11640, 11665--11670; GPR §2.2, eq.~(2.2.11), ll.~731--743}. GPR use the
dimensionless momenta $\sqrt{\ap/2}\,p_{L,R}$ and the dimensionless radius $R/\sqrt{\ap}$; in
those variables the map reads $R\mapsto1/R$. The proposition is a statement of elementary algebra,
and it is the part of this chapter that has been checked in Lean (\cref{sec:tduality:lean}).

Two features of \eqref{eq:tduality:pflip} deserve emphasis, because they are stronger than
the invariance of $M^2$.
\begin{itemize}[leftmargin=*]
\item \emph{The two chiral halves are treated differently.} The left-movers do not notice the
 duality at all, while the right-moving momentum is reflected. We shall see in
 \cref{sec:tduality:Y} that the same is true of the whole right-moving field, oscillators
 included. T-duality is a parity transformation acting on one chirality of the worldsheet
 theory\index{T-duality!as one-sided parity}.
\item \emph{The two Kaluza--Klein gauge fields are exchanged.} In \cref{ch:circle} we saw
 that the charge under the gauge field $A_\mu$ coming from the metric is proportional to
 $p_L+p_R=2n/R$, and the charge under $\tilde A_\mu$ coming from the $B$-field is proportional to
 $p_L-p_R=2wR/\ap$ \source{Tong §8.2.2, ll.~11527--11560}. Since $\mathsf T$ exchanges
 $p_L+p_R\leftrightarrow p_L-p_R$, it exchanges $A_\mu\leftrightarrow\tilde A_\mu$: what one
 observer calls ``the photon of geometry'' the dual observer calls ``the photon of the
 $B$-field''.
\end{itemize}

\begin{physicsbox}[Physical picture: what it means to move and what it means to wind]
Take $R\gg\ell_s$. Winding states cost an energy $\sim R/\ap$ and are out of reach; momentum
states are spaced by $1/R$ and, as $R\to\infty$, merge into a continuum. A continuum of
momentum states is precisely how a low-energy observer recognises a large dimension. Now take
$R\ll\ell_s$. The momentum states are heavy and drop out, but the winding states are now spaced
by $R/\ap=1/\tilde R$ and form a continuum as $R\to0$: the spectrum looks exactly as if a
\emph{new} large dimension of radius $\tilde R=\ap/R$ had opened up
\source{Tong §8.3, ll.~11643--11654}. A point particle has no winding states, which is why
Kaluza--Klein field theory on a small circle simply loses a dimension, while string theory
does not.
\end{physicsbox}

\begin{figure}[t]
\centering
\begin{tikzpicture}[x=2.6cm,y=1.35cm,>=Stealth]
  \draw[->] (0,0)--(3.45,0) node[right]{$R/\ell_s$};
  \draw[->] (0,0)--(0,3.4) node[left]{$M\ell_s$};
  \foreach \x in {1,2,3} \draw (\x,0.06)--(\x,-0.06) node[below]{\small$\x$};
  \foreach \y in {1,2,3} \draw (0.04,\y)--(-0.04,\y) node[left]{\small$\y$};
  \draw[thick,tierL,domain=0.31:3.2,samples=60] plot (\x,{1/\x});
  \draw[thick,tierC,domain=0:3.2] plot (\x,{\x});
  \draw[thick,tierL,dashed,domain=0.62:3.2,samples=60] plot (\x,{2/\x});
  \draw[thick,tierC,dashed,domain=0:1.62] plot (\x,{2*\x});
  \draw[dotted] (1,0)--(1,3.3);
  \fill (1,1) circle (1.6pt);
  \fill (1,2) circle (1.6pt);
  \node[tierL,below] at (2.9,0.33) {\small$n=1$};
  \node[tierL,above] at (2.9,0.72) {\small$n=2$};
  \node[tierC,right] at (2.85,2.6) {\small$w=1$};
  \node[tierC,right] at (1.55,3.0) {\small$w=2$};
  \node[above right] at (1,-0.02) {\small self-dual};
\end{tikzpicture}
\caption{The zero-mode contributions $|n|/R$ (momentum) and $|w|R/\ap$ (winding) to the mass, in
string units. The picture is symmetric under reflection $R/\ell_s\mapsto\ell_s/R$ together
with the exchange of the two families of curves. Light momentum states signal a large circle;
light winding states signal a large \emph{dual} circle. The full $M^2$ of
\eqref{eq:tduality:M2} adds the oscillator term, which does not depend on $R$.}
\label{fig:tduality:spectrum}
\end{figure}

\subsection{A worked example: $R=2\ell_s$ against $R=\ell_s/2$}\label{sec:tduality:example}

\begin{example}\label{ex:tduality:table}
Set $\ap=1$ and compare the circle of radius $R=2$ with its dual of radius $\tilde R=\frac12$.
At $R=2$, \eqref{eq:tduality:M2} reads $M^2=\frac14n^2+4w^2+2(N+\tilde N-2)$ with
$N-\tilde N=nw$; at $\tilde R=\frac12$ it reads $M^2=4\tilde n^2+\frac14\tilde w^2+2(N+\tilde N-2)$.
Listing all label sets $(n,w;N,\tilde N)$ with $M^2\le\frac94$ gives \cref{tab:tduality:example}.

\begin{table}[ht]
\centering
\caption{Lightest sectors of the bosonic string on a circle, $\ap=1$. ``Sectors'' counts label
sets $(n,w;N,\tilde N)$, not states: a sector with $N=1$ contains one state for each
polarization of the oscillator. The two columns are related by $n\leftrightarrow w$.}
\label{tab:tduality:example}
\small
\begin{tabular}{@{}rcp{0.36\textwidth}p{0.36\textwidth}@{}}\toprule
$M^2$ & sectors & $(n,w;N,\tilde N)$ at $R=2$ & $(\tilde n,\tilde w;N,\tilde N)$ at $\tilde R=\frac12$\\\midrule
$-4$ & 1 & $(0,0;0,0)$ & $(0,0;0,0)$\\
$-\frac{15}4$ & 2 & $(\pm1,0;0,0)$ & $(0,\pm1;0,0)$\\
$-3$ & 2 & $(\pm2,0;0,0)$ & $(0,\pm2;0,0)$\\
$-\frac74$ & 2 & $(\pm3,0;0,0)$ & $(0,\pm3;0,0)$\\
$0$ & 5 & $(\pm4,0;0,0)$, $(0,\pm1;0,0)$, $(0,0;1,1)$ & $(0,\pm4;0,0)$, $(\pm1,0;0,0)$, $(0,0;1,1)$\\
$\frac14$ & 2 & $(\pm1,0;1,1)$ & $(0,\pm1;1,1)$\\
$1$ & 2 & $(\pm2,0;1,1)$ & $(0,\pm2;1,1)$\\
$\frac94$ & 8 & $(\pm5,0;0,0)$, $(\pm3,0;1,1)$, $\pm(1,1;1,0)$, $\pm(1,-1;0,1)$ &
  $(0,\pm5;0,0)$, $(0,\pm3;1,1)$, $\pm(1,1;1,0)$, $\pm(-1,1;0,1)$\\\bottomrule
\end{tabular}
\end{table}

Three lessons can be read off. First, the two lists of masses and multiplicities coincide, as
\cref{thm:tduality:pLpR} guarantees; the table was generated by enumerating integers, not by
using the proposition, so it is an independent check. Second, a \emph{given} label does not
keep its mass: the state $(n,w)=(1,0)$ with $N=\tilde N=0$ has $M^2=-\frac{15}4$ at $R=2$ but
$M^2=0$ at $R=\frac12$. T-duality is not the statement that each state is insensitive to $R$;
it is the statement that the two \emph{collections} of states agree after relabelling. Third,
the bosonic string has tachyons ($M^2<0$) in both descriptions. This disease of the bosonic
string is unrelated to the circle and is cured in the superstring (\cref{ch:ghosts});
T-duality maps tachyons to tachyons.
\end{example}

\section{Four meanings of ``T-duality is a symmetry''}\label{sec:tduality:levels}

\Cref{thm:tduality:pLpR} is an identity between two lists of numbers. It is natural to ask how
much more is true. We distinguish four levels.
\begin{description}[leftmargin=1.2em]
\item[Level 1: the free spectrum.] The sets of masses, with multiplicities and charges, agree
 at radius $R$ and $\ap/R$. Equivalently the one-loop (torus) partition function agrees, because
 it is a trace over the spectrum. This is \cref{thm:tduality:pLpR}, extended from the zero modes
 to all states in \cref{sec:tduality:Y}.
\item[Level 2: the conformal field theory.] There is a one-to-one map between the local
 operators of the two worldsheet theories under which \emph{all} operator product expansions
 agree. Since tree-level string scattering amplitudes are correlation functions of vertex
 operators on the sphere (\cref{ch:ghosts}), level~2 implies that tree-level interactions
 agree. Two theories can have the same spectrum without being equivalent, in the same way that
 two drums of different shapes can sound the same notes, so level~2 is a genuinely stronger claim.
\item[Level 3: string perturbation theory.] The sum over worldsheets of all genera agrees.
 This requires a new ingredient: the string coupling must be transformed together with $R$
 (\cref{sec:tduality:dilaton}).
\item[Level 4: the exact theory.] The duality holds beyond perturbation theory. For the bosonic
 string this has no sharp meaning, since its perturbation series is already ill-defined because
 of the tachyon \source{GPR §2.3, ll.~1023--1029}. The best available argument, reviewed in
 \cref{sec:tduality:selfdual}, is that T-duality is a remnant of a \emph{gauge} symmetry, and gauge
 symmetries cannot be broken by quantum corrections.
\end{description}
Levels~1--3 are results of the literature, \TL. Only the zero-mode algebra underlying level~1
is kernel-checked in this book's library, \TA. No level is a statement about experiments: a
symmetry of string theory on a circle is a mathematical property of a model.

\section{The dual coordinate}\label{sec:tduality:Y}

\subsection{Flipping the right-movers}

The rule $p_L\mapsto p_L$, $p_R\mapsto-p_R$ suggests defining a new worldsheet field in which the
whole right-moving half of $X$ is reflected:
\begin{equation}\label{eq:tduality:Y}
  Y(\tau,\sigma)=X_L(\sigma^+)-X_R(\sigma^-)
\end{equation}
\source{Tong §8.3, ll.~11665--11678}\index{dual coordinate}. Let us find out what kind of field
$Y$ is. Insert the expansions \eqref{eq:tduality:modes}:
\begin{align}
  Y&=y+\tfrac12\ap(p_L-p_R)\,\tau+\tfrac12\ap(p_L+p_R)\,\sigma+\text{oscillators}\notag\\
   &=y+wR\,\tau+\ap\frac nR\,\sigma+\dots
    \;=\;y+\ap\frac{w}{\tilde R}\,\tau+n\tilde R\,\sigma+\dots,\qquad\tilde R=\frac{\ap}{R}.
    \label{eq:tduality:Ymodes}
\end{align}
Compare with $X=x+\ap\frac nR\tau+wR\sigma+\dots$. The field $Y$ is a free boson that advances by
$2\pi n\tilde R$ when $\sigma\to\sigma+2\pi$, and whose centre of mass carries momentum
$w/\tilde R$. In other words, $Y$ is a coordinate on a circle of radius $\tilde R=\ap/R$, on which
the string has winding number $n$ and momentum number $w$. The oscillators of $Y$ are
$\tilde\alpha_k$ and $-\alpha_k$, which is the rule
\begin{equation}\label{eq:tduality:osc}
  \alpha_k\longmapsto-\alpha_k,\qquad\tilde\alpha_k\longmapsto\tilde\alpha_k
\end{equation}
of \source{GPR §2.2, eq.~(2.2.12), ll.~742--750}. The level operators
$N=\sum_{k>0}\alpha_{-k}\alpha_k$ and $\tilde N$ are quadratic in the oscillators and therefore
invariant, which justifies the last entry of \eqref{eq:tduality:map}. So are the contributions
$\frac{\ap}4p_R^2+N$ and $\frac{\ap}4p_L^2+\tilde N$ of the compact direction to $L_0$ and
$\tilde L_0$, and more generally all the Virasoro generators, which are quadratic in
$\partial_\pm X$.

The relation between $X$ and $Y$ can be written without splitting into movers. From
$\partial_\tau=\partial_++\partial_-$ and $\partial_\sigma=\partial_+-\partial_-$,
\[
 \partial_\tau X=X_L'+X_R'=\partial_\sigma Y,\qquad
 \partial_\sigma X=X_L'-X_R'=\partial_\tau Y,
\]
or, covariantly,
\begin{equation}\label{eq:tduality:eps}
  \partial_\alpha X=\epsilon_{\alpha\beta}\,\partial^\beta Y,\qquad
  \epsilon_{\tau\sigma}=-\epsilon_{\sigma\tau}=+1,
\end{equation}
where the index is raised with the worldsheet Minkowski metric \source{Tong §8.3,
eq.~(8.9), ll.~11679--11690}. This is a two-dimensional version of electric--magnetic duality: the
one-form $\dd X$ is exchanged with its Hodge dual. It exchanges the equation of motion
$\partial_\alpha\partial^\alpha X=0$ with the Bianchi identity
$\epsilon^{\alpha\beta}\partial_\alpha\partial_\beta Y=0$, and the Noether charge of translations
(momentum, $\oint\dd\sigma\,\partial_\tau X$) with the topological charge (winding,
$\oint\dd\sigma\,\partial_\sigma X$). Note that \eqref{eq:tduality:eps} is non-local: $Y$ is
obtained from $X$ by integrating its derivatives, not by a pointwise change of variables. No
redefinition of the target-space coordinate $X\mapsto f(X)$ produces $Y$.

\subsection{From the spectrum to the operator algebra}

Pass to the Euclidean worldsheet and complex coordinates as in \cref{ch:cft}, where
left-movers become holomorphic and right-movers antiholomorphic: $X(z,\bar z)=X(z)+\bar X(\bar z)$.
The basic operator products are
\begin{equation}\label{eq:tduality:ope}
  X(z)X(w)\sim-\frac{\ap}2\ln(z-w),\qquad
  \bar X(\bar z)\bar X(\bar w)\sim-\frac{\ap}2\ln(\bar z-\bar w),\qquad
  X(z)\bar X(\bar w)\sim0 .
\end{equation}
Each is quadratic in the field to which the reflection is applied, so the fields
$Y(z)=X(z)$ and $\bar Y(\bar z)=-\bar X(\bar z)$ obey the same operator products, and the stress
tensors $T=-\frac1{\ap}\partial X\partial X$, $\bar T=-\frac1{\ap}\bar\partial\bar X\bar\partial\bar X$
are literally unchanged. Tong summarises: all OPEs of $Y$ coincide with those of $X$, and
this is sufficient to ensure that all interactions defined in the conformal field theory are
the same \source{Tong §8.3, ll.~11675--11678}.

It is instructive to see how this works for the operators that carry momentum and winding. A
state with labels $(n,w)$ is created by the vertex operator
\begin{equation}\label{eq:tduality:vertex}
  V_{n,w}(z,\bar z)={:}\exp\!\big(\ii p_LX(z)+\ii p_R\bar X(\bar z)\big){:}
   \;=\;{:}\exp\!\big(\ii p_LY(z)+\ii(-p_R)\bar Y(\bar z)\big){:}
\end{equation}
\source{Tong §8.2.2, ll.~11536--11545}. The second form is the key point: the \emph{same}
operator, expressed through $Y$, is a vertex operator with momenta $(p_L,-p_R)$, that is, with
labels $(w,n)$ on the circle of radius $\tilde R$. Nothing has been transformed; an operator
algebra has been given a second set of names. With the rule of \cref{ch:cft} that
$\ee^{\ii kX(z)}$ has weight $\ap k^2/4$, the weights of $V_{n,w}$ and the product of two such
operators are
\begin{gather}
  (h,\bar h)=\Big(\frac{\ap p_L^2}4,\frac{\ap p_R^2}4\Big),\qquad h-\bar h=nw\in\Z,
   \label{eq:tduality:weights}\\
  V_{n,w}(z,\bar z)\,V_{n',w'}(0)\sim
   z^{\ap p_Lp_L'/2}\;\bar z^{\ap p_Rp_R'/2}\;V_{n+n',w+w'}(0)+\dots,\notag\\
   \frac{\ap}2\big(p_Lp_L'-p_Rp_R'\big)=nw'+wn' .\label{eq:tduality:VV}
\end{gather}
The exponents depend on the right-moving momenta only through the product $p_Rp_R'$, which is
invariant under the simultaneous reflection of both factors. So the operator product
coefficients agree in the two descriptions: this is level~2. The last identity in
\eqref{eq:tduality:VV}, which follows by expanding \eqref{eq:tduality:pLpR}, shows that the
difference of the two exponents is an integer, so the product is single-valued when $z$
circles the origin. That integer, $nw'+wn'$, is a bilinear form on the charge lattice
$\Z^2\ni(n,w)$ which is independent of $R$ and invariant under $n\leftrightarrow w$; it will
reappear in \cref{ch:narain} as the form $\eta$ of signature $(1,1)$ whose automorphisms make
up the duality group $\Odd{1}$ (\cref{ch:odd}).

\begin{remark}
A complete treatment must also attach sign factors (``cocycles'') to the operators
$V_{n,w}$ so that they commute rather than anticommute when $nw'+wn'$ is odd, and must check
that the duality map can be chosen compatibly with them. This point is standard but is not
discussed in the sources of this book's library, and we do not pursue it.
\end{remark}

\begin{pitfallbox}[Common pitfall: ``the physics does not depend on $R$'']
T-duality does \emph{not} say that physics is independent of the radius. The spectrum at
$R=2\ell_s$ differs from the spectrum at $R=3\ell_s$; an observer can measure the spacing of the
light tower and distinguish them. What cannot be distinguished is $R=2\ell_s$ from
$R=\ell_s/2$. The symmetry is a discrete identification of the parameter space, not the
absence of a parameter.
\end{pitfallbox}

\section{The path integral, the string coupling and the dilaton}\label{sec:tduality:dilaton}

\subsection{A derivation that does not split into movers}

The argument of \cref{sec:tduality:Y} used the special solvability of the free boson. There is a
second derivation which works directly in the path integral, and which is the prototype of
the Buscher procedure\index{Buscher rules} of \cref{ch:buscher}, where it is applied to curved
backgrounds. We follow \source{Tong §8.3.1, ll.~11716--11803}; the general version is
\source{AAL §2.1, ll.~215--300 and 336--376}.

Write $X=R\varphi$ with $\varphi\sim\varphi+2\pi$, so that the radius appears as a coupling
constant. On a Euclidean worldsheet,
\begin{equation}\label{eq:tduality:Sphi}
  S[\varphi]=\frac{R^2}{4\pi\ap}\int\dd^2\sigma\;\partial_\alpha\varphi\,\partial^\alpha\varphi .
\end{equation}
\emph{Step 1: gauge the shift symmetry.} The action is invariant under
$\varphi\to\varphi+\lambda$. Introduce a worldsheet gauge field $A_\alpha\to A_\alpha-\partial_\alpha\lambda$
and the covariant derivative $D_\alpha\varphi=\partial_\alpha\varphi+A_\alpha$. By itself this
would change the theory, so add a Lagrange multiplier $\theta$ that forces the gauge field to
be trivial:
\begin{equation}\label{eq:tduality:Sgauged}
  S[\varphi,\theta,A]=\frac{R^2}{4\pi\ap}\int\dd^2\sigma\;D_\alpha\varphi\,D^\alpha\varphi
    +\frac{\ii}{2\pi}\int\dd^2\sigma\;\theta\,\epsilon^{\alpha\beta}\partial_\alpha A_\beta .
\end{equation}
\emph{Step 2: integrate out $\theta$.} This sets $\epsilon^{\alpha\beta}\partial_\alpha A_\beta=0$.
On a worldsheet with the topology of the plane, $A$ is then pure gauge and can be set to
zero, giving back \eqref{eq:tduality:Sphi}. On a worldsheet with handles a flat gauge field
can still have holonomies around the non-contractible cycles. If $\theta$ is itself periodic
with period $2\pi$, the sum over its winding sectors restricts these holonomies to multiples
of $2\pi$, and those are removed by gauge transformations $\lambda$ that wind around the cycle,
which are allowed because $\varphi$ is periodic. So \eqref{eq:tduality:Sgauged}, divided
by the volume of the gauge group, defines the same theory as \eqref{eq:tduality:Sphi} on every
worldsheet.

\emph{Step 3: integrate in the other order.} Fix the gauge $\varphi=0$ instead and
integrate the multiplier term by parts. The action becomes quadratic in $A$ with no
derivatives:
\[
  S=\frac{R^2}{4\pi\ap}\int\dd^2\sigma\,A_\alpha A^\alpha
    -\frac{\ii}{2\pi}\int\dd^2\sigma\,\epsilon^{\alpha\beta}(\partial_\alpha\theta)A_\beta .
\]
The Gaussian integral is done by solving the equation of motion of $A$,
\[
  \frac{R^2}{2\pi\ap}A^\beta=\frac{\ii}{2\pi}\epsilon^{\alpha\beta}\partial_\alpha\theta
  \quad\Longrightarrow\quad
  A^\beta=\ii\,\frac{\ap}{R^2}\,\epsilon^{\alpha\beta}\partial_\alpha\theta ,
\]
and substituting back. Using $\epsilon^{\alpha\beta}\epsilon_{\gamma\beta}=\delta^\alpha_\gamma$ in
Euclidean signature, the first term gives $-\frac{\ap}{4\pi R^2}(\partial\theta)^2$ and the
second $+\frac{\ap}{2\pi R^2}(\partial\theta)^2$, so
\begin{equation}\label{eq:tduality:Stheta}
  S[\theta]=\frac{\ap}{4\pi R^2}\int\dd^2\sigma\;\partial_\alpha\theta\,\partial^\alpha\theta
   =\frac{\tilde R^2}{4\pi\ap}\int\dd^2\sigma\;\partial_\alpha\theta\,\partial^\alpha\theta,
   \qquad\tilde R=\frac{\ap}{R}.
\end{equation}
This is \eqref{eq:tduality:Sphi} for a $2\pi$-periodic boson on a circle of radius $\tilde R$.
The equation of motion of $A$, at $\varphi=0$ where $A_\alpha$ stands for $D_\alpha\varphi$, is
the Euclidean form of \eqref{eq:tduality:eps} with $Y=\tilde R\theta$: the factor $\ii$ is the
usual consequence of the Wick rotation. The derivation never refers to the genus of the
worldsheet, so it supports level~3 of \cref{sec:tduality:levels}. One thing was discarded on
the way: the Gaussian integral over $A$ produces a field-independent factor, proportional to
$\sqrt{\ap}/R$ for each mode, which Tong sets aside with the remark that a more careful
treatment turns it into a shift of the dilaton \source{Tong §8.3.1, ll.~11795--11803}.
We now obtain that shift in two independent ways.

\subsection{The dilaton shift from the effective action}

Recall from \cref{ch:backgrounds} that the string coupling is the expectation value of the
dilaton, $g_s=\ee^{\Phi}$\index{string coupling}, and that a worldsheet of genus $g$ is weighted by $g_s^{2g-2}$.
The low-energy effective action in 26 dimensions is multiplied by $1/g_s^2$. Reducing on
the circle, the integral over $X^{25}$ gives a factor $2\pi R$, so the 25-dimensional
gravitational action comes with the coefficient
\begin{equation}\label{eq:tduality:G25}
  \frac{2\pi R}{\ell_s^{24}\,g_s^2}
\end{equation}
up to a number. An observer in 25 dimensions measures this combination, which is an inverse
Newton constant, not $R$ and $g_s$ separately. If the two descriptions are to be
indistinguishable we need $R/g_s^2=\tilde R/\tilde g_s^2$, that is,
\begin{equation}\label{eq:tduality:gs}
  \tilde g_s=g_s\sqrt{\frac{\tilde R}{R}}=\frac{\sqrt{\ap}}{R}\,g_s,\qquad
  \tilde\Phi=\Phi-\ln\frac{R}{\sqrt{\ap}}
\end{equation}
\source{Tong §8.3, eq.~(8.10), ll.~11691--11715}\index{dilaton!shift under T-duality}. This is a
plausibility argument, as Tong says. It can be restated by saying that the 25-dimensional dilaton
$\Phi_{25}=\Phi-\frac12\ln(R/\sqrt{\ap})$ is T-duality invariant: indeed
$\tilde\Phi-\frac12\ln(\tilde R/\sqrt{\ap})=\Phi-\ln(R/\sqrt{\ap})+\frac12\ln(R/\sqrt{\ap})$. The
same rule, in the form $\tilde\phi=\phi-\frac12\ln g_{00}$ with $g_{00}=R^2/\ap$ the metric
component along the circle, comes out of a regularized treatment of the Gaussian determinant
\source{AAL §2.1, eqs.~(2.1.6)--(2.1.7), ll.~296--312}; it is one of the Buscher rules of
\cref{ch:buscher}. Two consistency checks: at $R=\sqrt{\ap}$ the coupling does not change, as it
must at a fixed point; and applying \eqref{eq:tduality:gs} twice returns $g_s$, since
$(\sqrt{\ap}/\tilde R)(\sqrt{\ap}/R)=1$.

\subsection{The dilaton shift from the sum over genera}

Giveon, Porrati and Rabinovici obtain the same result from the partition function. Let $r=R/\sqrt{\ap}$
be the dimensionless radius. On a worldsheet of genus $g$ the compact boson has $2g$
integers, its windings around the $a$- and $b$-cycles. After a Poisson resummation over half of
them, the sum over these sectors becomes a sum over $g$ copies of the pairs $(p_L,p_R)$ that is
manifestly invariant under $r\to1/r$ with the integers exchanged, multiplied by an explicit
power of the radius:
\begin{equation}\label{eq:tduality:Zg}
  Z_g(r)=r^{1-g}\,F_g\big(r\big),\quad F_g(1/r)=F_g(r)\qquad\Longrightarrow\qquad
  Z_g(1/r)=r^{2g-2}\,Z_g(r)
\end{equation}
\source{GPR §2.3, eqs.~(2.3.6)--(2.3.9), ll.~930--997}. One power of $r$ comes from the volume of
the zero mode of $X$, and $r^{-g}$ from the resummation. At genus one the factor is absent:
$Z_1(1/r)=Z_1(r)$. This is level~1, since the torus amplitude is a trace over the spectrum
and the integers $n,w$ are summation variables \source{GPR §2.2, ll.~778--786}. At genus
$g\ne1$ the partition functions differ, but only by a power of $r$ that depends on the
genus exactly as the weight $g_s^{2g-2}$ does. Hence
\begin{equation}\label{eq:tduality:Zfull}
  \sum_{g\ge0}\tilde g_s^{\,2g-2}Z_g(1/r)=\sum_{g\ge0}\big(\tilde g_s\,r\big)^{2g-2}Z_g(r)
  =\sum_{g\ge0}g_s^{2g-2}Z_g(r)\qquad\text{iff}\qquad\tilde g_s=\frac{g_s}{r},
\end{equation}
which is \eqref{eq:tduality:gs} again. GPR write the weight as $\ee^{(1-g)\Phi_{\rm GPR}}$, so that
$\Phi_{\rm GPR}=-2\Phi$, and state the result as $\Phi_{\rm GPR}\to\Phi_{\rm GPR}+2\ln r$
\source{GPR §2.3, eqs.~(2.3.10)--(2.3.11), ll.~998--1029}; converting, $\tilde\Phi=\Phi-\ln r$, in
agreement with \eqref{eq:tduality:gs}. Their conclusion is worth quoting in substance: a
compactification of small radius is equivalent to one of large radius \emph{provided} the
dilaton expectation value is changed as well.

\begin{pitfallbox}[Common pitfall: forgetting the coupling]
``$R\to\ap/R$ is a symmetry of string theory'' is false as it stands: by
\eqref{eq:tduality:Zg} only the torus amplitude is invariant under the inversion of $R$ alone. The
symmetry is $(R,g_s)\to(\ap/R,\;g_s\sqrt{\ap}/R)$. For $R\ll\ell_s$ the dual coupling is
\emph{larger} than the original one: a weakly coupled string on a very small circle is dual to
a string on a large circle whose ten- or twenty-six-dimensional coupling is stronger, although
the lower-dimensional coupling \eqref{eq:tduality:G25} is the same.
\end{pitfallbox}

\section{The self-dual radius}\label{sec:tduality:selfdual}

\subsection{The fixed point and the moduli space}

The map $R\mapsto\ap/R$ on the positive reals has exactly one fixed point,
\begin{equation}\label{eq:tduality:sd}
  R=\frac{\ap}{R}\iff R^2=\ap\iff R=\sqrt{\ap}=\ell_s ,
\end{equation}
the \emph{self-dual radius}\index{self-dual radius}. The radius is not a fixed parameter of the
theory but the expectation value of a massless scalar field, the radion $\sigma$ of
\cref{ch:circle}; such a parameter is called a \emph{modulus}\index{moduli space!of the circle},
and the set of its inequivalent values is the moduli space. Naively this is the half-line
$R\in(0,\infty)$. T-duality tells us that $R$ and $\ap/R$ label the same theory, so the true
moduli space is the quotient
\begin{equation}\label{eq:tduality:moduli}
  \mathcal M_{S^1}=(0,\infty)\big/\big(R\sim\ap/R\big)\;\cong\;[\sqrt{\ap},\infty).
\end{equation}
In the logarithmic coordinate $x=\ln(R/\sqrt{\ap})$ the duality is the reflection
$x\mapsto-x$ of the real line, and $\mathcal M_{S^1}=\R/\Z_2=[0,\infty)$: a half-line whose
end point is the self-dual radius (\cref{fig:tduality:moduli}). Every circle compactification is
equivalent to one with $R\ge\ell_s$. The coordinate $x$ is also the natural one for measuring
distances in moduli space, a fact that becomes important in \cref{ch:swampland}.

\begin{figure}[ht]
\centering
\begin{tikzpicture}[>=Stealth]
  \draw[->] (-4.6,0)--(4.8,0) node[right]{$x=\ln(R/\ell_s)$};
  \fill (0,0) circle (2pt) node[below=3pt]{\small$0$};
  \draw[<->,tierL,thick] (-2.5,0.2) .. controls (-1.2,1.0) and (1.2,1.0) .. (2.5,0.2);
  \node[tierL] at (0,1.05) {\small$\mathsf T:\;x\mapsto-x$};
  \fill[tierC] (2.5,0) circle (1.6pt) node[below=3pt]{\small$R$};
  \fill[tierC] (-2.5,0) circle (1.6pt) node[below=3pt]{\small$\ap/R$};
  \node[below] at (-4.1,-0.05) {\small$R\to0$};
  \node[below] at (4.2,-0.05) {\small$R\to\infty$};
  \draw[->] (0,-0.75)--(0,-1.45) node[midway,right]{\small quotient by $\Z_2$};
  \draw[very thick] (0,-2)--(4.4,-2);
  \draw[very thick,dotted] (4.4,-2)--(5.0,-2);
  \fill (0,-2) circle (2.4pt);
  \node[below] at (0,-2.1) {\small$R=\ell_s$: $SU(2)\times SU(2)$};
  \node[above] at (2.9,-1.95) {\small$U(1)\times U(1)$};
  \node[below] at (4.5,-2.1) {\small$R\to\infty$};
\end{tikzpicture}
\caption{Top: T-duality acts on the logarithm of the radius as a reflection. Bottom: the
moduli space of the circle compactification is the quotient, a half-line. At its end point
the gauge symmetry is enhanced.}
\label{fig:tduality:moduli}
\end{figure}

\subsection{Extra massless states}

End points of moduli spaces are where something happens. Set $R=\sqrt{\ap}$ in
\eqref{eq:tduality:M2}:
\begin{equation}\label{eq:tduality:M2sd}
  \ap M^2=n^2+w^2+2(N+\tilde N)-4,\qquad N-\tilde N=nw .
\end{equation}
We look for all solutions with $M^2=0$. Since $N+\tilde N\ge|N-\tilde N|=|nw|$, we need
$n^2+w^2+2|nw|\le4$, that is $(|n|+|w|)^2\le4$. The possibilities are therefore:
\begin{enumerate}[leftmargin=*,label=(\roman*)]
\item $n=w=0$, $N=\tilde N=1$: the states
 $\alpha^M_{-1}\tilde\alpha^{M'}_{-1}\ket{0;p}$ that are massless at every radius. With both indices
 non-compact they give the graviton, the $B$-field and the dilaton in 25 dimensions; with one
 index along the circle, the two gauge fields $A_\mu$, $\tilde A_\mu$; with both along the
 circle, the radion \source{Tong §8.2.1, ll.~11491--11526}.
\item $(n,w)=(\pm2,0)$ or $(0,\pm2)$, $N=\tilde N=0$: four scalars, momentum and winding modes of
 the tachyon.
\item $(n,w)=\pm(1,1)$, so $nw=1$ and $(N,\tilde N)=(1,0)$: the states $\alpha^M_{-1}\ket{n,w}$. For
 $M=\mu$ these are two \emph{vectors}; for $M=25$, two scalars.
\item $(n,w)=\pm(1,-1)$, so $nw=-1$ and $(N,\tilde N)=(0,1)$: the states
 $\tilde\alpha^M_{-1}\ket{n,w}$; again two vectors and two scalars.
\end{enumerate}
This agrees with the list of \source{Tong §8.2.3, ll.~11561--11612} (where the levels of the
last item should read $N=0$, $\tilde N=1$, as level matching requires) and with the direct
enumeration of \cref{ex:tduality:table} repeated at $R=1$, which finds nine label sets at
$M^2=0$. Away from the self-dual radius the new states are massive. For the vectors of
item~(iii),
\begin{equation}\label{eq:tduality:W}
  M^2=\frac1{R^2}+\frac{R^2}{\ap^2}-\frac2{\ap}=\Big(\frac1R-\frac R{\ap}\Big)^2=p_R^2,
  \qquad M\simeq\frac{2}{\ap}\,|R-\sqrt{\ap}|\ \text{ near }R=\sqrt{\ap}.
\end{equation}

\subsection{Enhanced gauge symmetry and T-duality as a gauge transformation}

Massless vectors that carry charge under $U(1)\times U(1)$ can only be consistent as the
gauge bosons of a non-abelian group. The charges are $(p_L+p_R,p_L-p_R)\propto(n,w)$, equal to
$\pm(1,1)$ and $\pm(1,-1)$; together with the two neutral photons this is the content of
\begin{equation}\label{eq:tduality:su2}
  U(1)\times U(1)\;\longrightarrow\;SU(2)\times SU(2)\qquad\text{at }R=\sqrt{\ap}
\end{equation}
\source{Tong §8.2.3, ll.~11608--11621}\index{enhanced gauge symmetry}. In terms of $p_L,p_R$: the
vectors of item~(iii) have $p_R=0$, $p_L=\pm2/\sqrt{\ap}$ and are charged only under the
left-moving combination of the photons; those of item~(iv) have $p_L=0$ and are charged only under
the right-moving one. So one $SU(2)$ is carried by left-movers and one by right-movers. The
scalars organise accordingly. There are $1+4+2+2=9$ of them, and they fill the representation
$(\mathbf3,\mathbf3)$; GPR describe them as scalars of spin $(1,1)$ under $SU(2)_L\times SU(2)_R$
\source{GPR §2.6.1, ll.~2416--2422}. Moving away from the self-dual radius is giving an
expectation value to one of these scalars, the radion, and \eqref{eq:tduality:W} is the Higgs
mechanism\index{Higgs mechanism}: the charged vectors acquire a mass proportional to the expectation value, and
$SU(2)\times SU(2)$ is broken to its Cartan subgroup $U(1)\times U(1)$.

On the worldsheet the enhancement is visible as extra conserved currents. At any radius
$\partial X(z)$ is a holomorphic current of weight $(1,0)$. By \eqref{eq:tduality:weights} the
operator $V_{n,w}$ has weights $(1,0)$ precisely when $p_R=0$ and $\ap p_L^2=4$, which at
$R=\sqrt{\ap}$ is solved by $(n,w)=\pm(1,1)$. So exactly at the self-dual radius there are three
holomorphic currents,
\begin{equation}\label{eq:tduality:currents}
  J^3(z)\propto\ii\,\partial X(z),\qquad J^\pm(z)={:}\exp\!\big(\pm2\ii X(z)/\sqrt{\ap}\big){:},
\end{equation}
and three antiholomorphic ones $\bar J^3,\bar J^\pm$ built in the same way from $\bar X$; they
generate the two $SU(2)$ current algebras \source{GPR §2.6.1, ll.~2474--2481}.

We can now see what T-duality is at the self-dual point. It acts as $\bar X\mapsto-\bar X$ and
trivially on $X(z)$, hence
\begin{equation}\label{eq:tduality:weyl}
  \bar J^3\longmapsto-\bar J^3,\qquad\bar J^\pm\longmapsto\bar J^\mp,\qquad J^a\longmapsto J^a .
\end{equation}
This is a rotation by $\pi$ about the $1$-axis of $SU(2)_R$: the non-trivial element of the Weyl
group\index{Weyl group} of $SU(2)_R$. It is therefore a \emph{gauge transformation} of the
25-dimensional theory. A small deformation $R=\sqrt{\ap}(1+\delta)$ is a perturbation of the
worldsheet action by the operator $\delta\cdot J^3\bar J^3$, which the rotation maps to
$-\delta\cdot J^3\bar J^3$, i.e.\ to $R=\sqrt{\ap}(1-\delta)\approx\ap/R$. The theories at $R$ and
at $\ap/R$ are thus related by a gauge transformation of the enhanced symmetry
\source{GPR §2.6.1, ll.~2436--2481}. This is the conceptual reason to believe in level~4 of
\cref{sec:tduality:levels}: two configurations related by a gauge transformation are not merely
similar but are the same physical state, and gauge redundancies are not broken by loop or
non-perturbative effects.

\begin{remark}
The bosonic mechanism of enhancement uses momentum and winding modes of the tachyon and
does not operate in the type~II superstrings; it does operate in the heterotic string
\source{Tong §8.2.3, ll.~11622--11625}. T-duality itself survives in all of them, with
modifications discussed in \cref{ch:dbranes}.
\end{remark}

\section{Minimal length as physics}\label{sec:tduality:minlength}

\subsection{What an observer can measure}

All three of our sources draw the same moral: T-duality is a mechanism by which string theory
exhibits a minimal length\index{minimal length}. As the circle shrinks below $R=\sqrt{\ap}$ ``the
theory acts as if the circle is growing again, with winding modes playing the role of momentum
modes'' \source{Tong §8.3, ll.~11659--11662}; a small radius ``is completely equivalent to a very
large one'' \source{GPR §2.2, ll.~784--793}. To turn the slogan into a statement we must say how a
length is measured. An observer in the non-compact dimensions has no ruler to lay along the
circle. What she can do is find the lightest tower of equally spaced states and \emph{define}
the radius as the inverse spacing, because this is how a Kaluza--Klein circle shows up in
experiments. For $R>\ell_s$ the lightest tower is the momentum tower, with spacing $1/R$; for
$R<\ell_s$ it is the winding tower, with spacing $R/\ap$. The operationally defined radius is
therefore
\begin{equation}\label{eq:tduality:Robs}
  R_{\rm obs}=\max\!\Big(R,\frac{\ap}{R}\Big)\;\ge\;\sqrt{\ap}.
\end{equation}
\'Alvarez, \'Alvarez-Gaum\'e and Lozano make the same point with correlation functions:
distances extracted from the large-separation behaviour of correlators come out differently
for pure momentum probes and for pure winding probes, the latter seeing a circle of radius
$\ap/R$, and this leads ``in a natural way'' to restricting the outcome of distance measurements
to the interval $d\in(1,\infty)$ in string units \source{AAL §6, ll.~1289--1301}. The
parameter $R$ can be made as small as we like; the \emph{measured} size cannot.

\begin{example}
Let $R=0.1\,\ell_s$. The momentum states have masses $10|n|/\ell_s$; the winding states have
masses $0.1|w|/\ell_s$. The light tower has spacing $0.1/\ell_s$, and the observer reports a circle
of radius $R_{\rm obs}=10\,\ell_s$. If she also measures the lower-dimensional Newton constant she
finds the value appropriate to a circle of radius $10\ell_s$ and coupling $\tilde g_s=10\,g_s$,
by \eqref{eq:tduality:gs}. There is no measurement whose outcome is ``$0.1\,\ell_s$''.
\end{example}

\subsection{The dual-scale function}

The function \eqref{eq:tduality:Robs} is duality invariant but not smooth at the self-dual
point. This programme uses instead the simplest smooth invariant built from the two radii,
the \emph{effective} or \emph{dual scale}\index{dual-scale principle}\index{effective scale}
\begin{equation}\label{eq:tduality:Reff}
  R_{\rm eff}(R)=R+\frac{\ap}{R}=2\sqrt{\ap}\,\cosh x,\qquad x=\ln\frac{R}{\sqrt{\ap}} .
\end{equation}
It treats $R$ and $\tilde R$ on an equal footing, and:
\begin{proposition}\label{thm:tduality:Reff}
For all $R>0$: (a) $R_{\rm eff}(\ap/R)=R_{\rm eff}(R)$; (b) $R_{\rm eff}(R)\ge2\sqrt{\ap}$, with
equality if and only if $R=\sqrt{\ap}$; (c) $R_{\rm obs}\le R_{\rm eff}\le2R_{\rm obs}$.
\end{proposition}
\begin{proof}
(a) is the symmetry of a sum. (b) follows from
\[
  R+\frac{\ap}{R}-2\sqrt{\ap}=\Big(\sqrt R-\sqrt{\ap/R}\Big)^2=\frac{(R-\sqrt{\ap})^2}{R}\ge0,
\]
which vanishes only for $R=\sqrt{\ap}$; this is the inequality between the arithmetic and the
geometric mean of $R$ and $\ap/R$, whose product is fixed. (c) A sum of two positive numbers
lies between the larger of them and twice the larger.
\end{proof}
Part~(c) says that $R_{\rm eff}$ agrees with the operational radius up to a factor of at most two
and coincides with it asymptotically, $R_{\rm eff}/R_{\rm obs}\to1$ as $R\to0$ or $\infty$. A
closely related invariant appears directly in the mass formula. With $r=R/\sqrt{\ap}$ the
zero-mode part of \eqref{eq:tduality:M2} is
\begin{equation}\label{eq:tduality:massform}
  \ap M^2_{\text{zero modes}}=\frac{n^2}{r^2}+w^2r^2=
   \begin{pmatrix}w&n\end{pmatrix}\begin{pmatrix}r^2&0\\0&r^{-2}\end{pmatrix}
   \begin{pmatrix}w\\n\end{pmatrix},
\end{equation}
a quadratic form in the doubled charge vector $Z=(w,n)$\index{doubled charge vector} whose matrix,
$\HH=\diag(r^2,r^{-2})$, is the $d=1$, $B=0$ generalized metric\index{generalized metric} of
\cref{ch:genmetric}. T-duality acts on it as $\HH\mapsto\eta\HH\eta$ with
$\eta=\big(\begin{smallmatrix}0&1\\1&0\end{smallmatrix}\big)$, which exchanges the two diagonal
entries. Its trace is invariant and equals
\begin{equation}\label{eq:tduality:trace}
  \Tr\HH=r^2+r^{-2}=\Big(r+\frac1r\Big)^2-2=\frac{R_{\rm eff}^2}{\ap}-2\;\ge\;2 .
\end{equation}
The bound $\Tr\HH\ge2$ is the $d=1$ case of the inequality $\Tr G+\Tr G^{-1}\ge2d$ for tori
proved in \cref{ch:dualscale}.

The status of these statements must be kept apart. The inequality of
\cref{thm:tduality:Reff}(b) and the identity \eqref{eq:tduality:trace} are mathematics, \TA{} in
the units $\ap=1$ (\cref{sec:tduality:lean}); the remaining parts of the proposition are proved
above by hand. That T-duality implies a minimal \emph{measurable} compactification
radius for closed-string probes is the view of the literature quoted above, \TL. That
$R_{\rm eff}$, rather than $R_{\rm obs}$ or any other invariant, is ``the'' physical scale, and
that it regulates cosmological singularities, is a proposal of this programme, \TC, examined
in \cref{ch:dualscale,ch:cosmology}. One caveat from the literature belongs here already:
the minimal length concerns closed strings as probes. Polchinski reports evidence that D-branes
probe length scales shorter than the string length, ``in conflict with the idea that the
string length is a minimum distance'' \source{Pol, ll.~3020--3025}. The minimal length of this
section is a property of perturbative closed-string physics on a circle, not a theorem about
quantum gravity.

\section{What the machine has checked}\label{sec:tduality:lean}

Four Lean files touch the material of this chapter. They differ greatly in how much they
say, and we go through them from the most faithful to the least. As everywhere in this book,
a kernel-checked statement is a statement about the Lean objects that appear in it; that those
objects model a string on a circle is our reading and is never \TA. According to
\lean{\#print axioms}, the theorems quoted below use no axioms beyond the standard
\lean{propext}, \lean{Classical.choice} and \lean{Quot.sound}.

\subsection{Zero-mode T-duality over the rationals}

File \path{StringTheoryFormalization/UseCases/TDualityMassSpectrum.lean}, namespace\\
\lean{StringTheory.UseCases.TDuality}. Units $\ap=1$; argument order $(n,w,R)$.

\begin{leanbox}[In Lean: the momenta and their transformation]
\begin{leancode}
def leftMomentum (n w R : ℚ) : ℚ := n / R + w * R

def rightMomentum (n w R : ℚ) : ℚ := n / R - w * R

def momentumMassSq (n w R : ℚ) : ℚ :=
  (leftMomentum n w R) ^ 2 + (rightMomentum n w R) ^ 2

theorem tduality_fixes_left_momentum (n w R : ℚ) (hR : R ≠ 0) :
    leftMomentum w n (1 / R) = leftMomentum n w R := by ...

theorem tduality_flips_right_momentum (n w R : ℚ) (hR : R ≠ 0) :
    rightMomentum w n (1 / R) = - rightMomentum n w R := by ...
\end{leancode}
\textbf{Reading.} The definitions are \eqref{eq:tduality:pLpR} at $\ap=1$, as functions
$\Q^3\to\Q$. Note that $n$, $w$ are arbitrary rationals here, not integers: the theorems hold for
all rational values, in particular the integral ones, and integrality is never used. The two
theorems are exactly \eqref{eq:tduality:pflip}: swapping the first two arguments and
replacing $R$ by $1/R$ fixes the first function and negates the second, for every $R\ne0$
(positivity of $R$ is not needed). \textbf{Proof idea.} Unfold, clear denominators with
\lean{field\_simp} using $R\ne0$, and close the polynomial identity with \lean{ring}.
\end{leanbox}

\begin{leanbox}[In Lean: invariance of the zero-mode mass]
\begin{leancode}
theorem tduality_invariant_mass_squared (n w R : ℚ) (hR : R ≠ 0) :
    momentumMassSq w n (1 / R) = momentumMassSq n w R := by ...
\end{leancode}
\textbf{Reading.} The quantity $p_L^2+p_R^2=2(n^2/R^2+w^2R^2)$ is invariant under the duality
map. By \eqref{eq:tduality:M2} this is \emph{twice} the zero-mode contribution to $M^2$; the
factor is harmless but should be remembered when comparing with other files. The oscillator
term $\frac2{\ap}(N+\tilde N-2)$ and level matching are not modelled: the theorem is the
zero-mode core of level~1 of \cref{sec:tduality:levels}, and says nothing about levels~2--4.
\textbf{Proof idea.} Rewrite with the two previous theorems; then $(-p_R)^2=p_R^2$ by
\lean{ring}. The proof thus follows the logic of \cref{thm:tduality:pLpR}: the stronger
statement about $p_L,p_R$ first, the mass as a corollary.
\end{leanbox}

\begin{leanbox}[In Lean: the self-dual radius]
\begin{leancode}
theorem self_dual_radius_unique (R : ℚ) (hR : 0 < R)
    (hFix : (1 : ℚ) / R = R) : R = 1 := by ...
\end{leancode}
\textbf{Reading.} A positive rational fixed by $R\mapsto1/R$ equals $1$. This is the
uniqueness half of \eqref{eq:tduality:sd} in units $\ap=1$; existence ($1/1=1$) is trivial and
not stated. \textbf{Proof idea.} \lean{field\_simp} turns the hypothesis into $1=R^2$ (up to
orientation); \lean{nlinarith}, given the hints $(R-1)^2\ge0$ and $(R+1)^2\ge0$ and $R>0$, excludes
the root $R=-1$.
\end{leanbox}

\begin{pitfallbox}[Common pitfall: units hidden in a number field]
Working over $\Q$ with $\ap=1$ is not an innocent choice. For general $\ap$ the fixed-point
equation is $R^2=\ap$, which has \emph{no} rational solution when, say, $\ap=2$. A version of
\lean{self\_dual\_radius\_unique} with a parameter $\ap$ must live over $\R$ and conclude
$R=\sqrt{\ap}$ (\cref{ex:tduality:leansqrt}). Setting $\ap=1$ is legitimate, since it is a
choice of unit of length, but then ``$R=1$'' means ``$R=\ell_s$''.
\end{pitfallbox}

The same momenta are consumed by a theorem of the newer library, \lean{massForm\_circle} in
the namespace \lean{DualScaleStream2.DFT}, which states that the quadratic form
\eqref{eq:tduality:massform}, evaluated with the generalized metric of $G=R^2$, $B=0$ on the
charge vector $(w,n)$, equals \lean{momentumMassSq n w R} divided by two. It is discussed in
\cref{ch:genmetric}; we mention it because it fixes the convention $Z=(w,n)$ and confirms the
factor of two noted above.

\subsection{The dual-scale bound over the reals}

File \path{DualScaleStream2/DualScale/TraceBound.lean}, with namespace\\
\lean{DualScaleStream2.DualScale}. Here \lean{dualScale G} is defined as the trace of the
generalized metric at $B=0$, for a real $d\times d$ matrix $G$; the file proves
$\lean{dualScale}\,G=\Tr G+\Tr G^{-1}$ (\lean{dualScale\_eq}), its invariance under
$G\mapsto G^{-1}$ (\lean{dualScale\_inv}) and the bound $2d$ for positive-definite $G$
(\lean{dualScale\_ge}), which belong to \cref{ch:dualscale}. Two statements concern the circle.

\begin{leanbox}[In Lean: the circle case of the dual scale]
\begin{leancode}
theorem dualScale_circle (R : ℝ) (hR : R ≠ 0) :
    dualScale (!![R ^ 2] : Matrix (Fin 1) (Fin 1) ℝ)
      = (R + R⁻¹) ^ 2 - 2 := by ...

theorem circle_effective_scale_ge_two (R : ℝ) (hR : 0 < R) :
    2 ≤ R + R⁻¹ := by ...
\end{leancode}
\textbf{Reading.} The first theorem is \eqref{eq:tduality:trace}: for the $1\times1$ metric
$G=(R^2)$ the trace of the generalized metric is $(R+1/R)^2-2$. The second is
\cref{thm:tduality:Reff}(b) at $\ap=1$, over the real numbers: $R+1/R\ge2$ for $R>0$. The
equality case and the invariance under $R\mapsto1/R$ are not part of these two statements
(the invariance is the $d=1$ case of \lean{dualScale\_inv}). Nothing in either statement
refers to strings, measurement or cosmology. \textbf{Proof idea.} For the first, compute
the inverse of a $1\times1$ matrix through the adjugate formula, then \lean{field\_simp} and
\lean{ring}. For the second, rewrite $R+R^{-1}-2=(R-1)^2/R$ and use that a square divided by a
positive number is non-negative: the proof of \cref{thm:tduality:Reff}(b), line by line.
\end{leanbox}

\subsection{Arithmetic shadows in the older libraries}

Two further files carry T-duality in their names. They belong to the older layer of the
project, in which physical quantities are replaced by integers or natural numbers, and their
theorems are far weaker than their documentation strings suggest. We record them so that the
reader who meets them in the repository knows what they are.

\begin{leanbox}[In Lean: \texttt{DualScaleValidation.UseCase1}]
\begin{leancode}
def alpha_prime : Nat := 1
def effective_dual_scale_numerator (R : Nat) : Nat := R * R + alpha_prime

theorem buscher_log_involution (x : Int) : -(-x) = x := by ...
theorem dual_scale_self_dual_value :
    effective_dual_scale_numerator 1 = 2 := by ...
theorem self_dual_is_global_minimum (R : Nat) (h : R ≥ 1) :
    effective_dual_scale_numerator R ≥ 2 := by ...
\end{leancode}
\textbf{Reading.} The first theorem is the fact that negation of integers is an involution;
it is an arithmetic shadow of ``$x\mapsto-x$ squares to the identity'' in
\cref{fig:tduality:moduli}, with the real coordinate $x$ replaced by an integer. The second
is $1\cdot1+1=2$. The third says $R^2+1\ge2$ for natural numbers $R\ge1$. Since
$R^2+\ap=R\cdot R_{\rm eff}(R)$, this is a statement about $R\cdot R_{\rm eff}$, not about
$R_{\rm eff}$: the inequality equivalent to $R_{\rm eff}\ge2$ would be $R^2+1\ge2R$, which is
not what is stated, and on natural numbers $R\ge1$ the dual region $R<1$ is absent altogether.
The file's comments speak of minimal length and of the elimination of cosmological
singularities; none of that is contained in these statements. A faithful version is
\lean{circle\_effective\_scale\_ge\_two} above. \textbf{Proof idea.} \lean{omega} and
\lean{rfl}.
\end{leanbox}

\begin{leanbox}[In Lean: \texttt{StringTheory.StringDynamics} (file \texttt{TDualityGysin.lean})]
\begin{leancode}
structure TDualState where
  momentumNum : ℤ   -- n (KK momentum)
  windingNum : ℤ    -- w (winding number)

def tDualityAction (s : TDualState) : TDualState :=
  { momentumNum := s.windingNum
    windingNum  := s.momentumNum }

theorem tduality_involution (s : TDualState) :
    tDualityAction (tDualityAction s) = s := by ...
\end{leancode}
\textbf{Reading.} Swapping the two components of a pair of integers twice gives back the
pair: a correct, if minimal, formal statement of $\mathsf T^2=1$ on the charges, with no
radius, mass or Hilbert space involved. The same file defines \lean{tDualityRadius} as the quotient $\ap/R$ of two reals
(its positivity hypotheses are unused) and \lean{gysinPushforward}, which evaluates an integer
sequence at $0$; despite its name the latter constructs no cohomology and no fibre
integration, as the file's own header states. \textbf{Proof idea.} \lean{simp} after unfolding.
\end{leanbox}

\subsection{Tier table and what is missing}

\begin{table}[ht]
\centering
\caption{Epistemic status of the main claims of this chapter.}\label{tab:tduality:tiers}
\small
\begin{tabular}{@{}p{0.47\textwidth}cp{0.40\textwidth}@{}}\toprule
Claim & Tier & Basis\\\midrule
$p_L\mapsto p_L$, $p_R\mapsto-p_R$ under $R\mapsto1/R$, $n\leftrightarrow w$ ($\ap=1$, over $\Q$) & \TA &
  \lean{tduality\_fixes\_left\_momentum}, \lean{tduality\_flips\_right\_momentum}\\
$p_L^2+p_R^2$ invariant & \TA & \lean{tduality\_invariant\_mass\_squared}\\
$R=1$ is the only positive rational fixed point of $R\mapsto1/R$ & \TA & \lean{self\_dual\_radius\_unique}\\
$\Tr\HH=(R+1/R)^2-2$; $R+1/R\ge2$ on $\R_{>0}$ & \TA &
  \lean{dualScale\_circle}, \lean{circle\_effective\_scale\_ge\_two}\\
$(n,w)\mapsto(w,n)$ is an involution of $\Z^2$ & \TA & \lean{tduality\_involution}\\
Full spectrum, oscillators included, agrees (level 1) & \TL & Tong §8.3; GPR §2.2\\
All OPEs agree; tree amplitudes agree (level 2) & \TL & Tong §8.3; AAL §2.1\\
All genera agree with $\tilde g_s=g_s\sqrt{\ap}/R$ (level 3) & \TL & GPR §2.3; Tong §8.3.1\\
$SU(2)\times SU(2)$ at $R=\sqrt{\ap}$; T-duality is a Weyl reflection & \TL & Tong §8.2.3; GPR §2.6.1\\
Minimal measurable radius for closed-string probes & \TL & Tong §8.3; GPR §2.2; AAL §6\\
Masses in \cref{tab:tduality:example}; nine massless label sets at $R=\ell_s$ & --- & enumeration
  by script; symbolic check, not Tier~A\\
$R_{\rm eff}$ is the physical scale; it regulates singularities & \TC & this programme\\\bottomrule
\end{tabular}
\end{table}

What would a faithful formalization of T-duality on a circle need? In increasing order of
difficulty: (i) the integrality of $n,w$ and the level-matching constraint, so that the
\emph{set} of states and its bijection can be stated, which is a modest extension of the present
file; (ii) the Fock space of the oscillators and a unitary operator implementing
\eqref{eq:tduality:osc}, intertwining the two Hamiltonians, which needs infinite-dimensional
Hilbert spaces but no new mathematics; (iii) the vertex operator algebra of the lattice
$\Z^2$ with the form $nw'+wn'$, and the statement that \eqref{eq:tduality:map} induces an isomorphism
of it, which is level~2 and would require a library of vertex algebras that Mathlib does not
currently provide; (iv) higher-genus partition functions and \eqref{eq:tduality:Zg}, which need
theta functions on Riemann surfaces. None of (ii)--(iv) exists in this project.

\begin{summarybox}
\begin{itemize}[leftmargin=*]
\item T-duality of the circle is the map $R\mapsto\ap/R$, $n\leftrightarrow w$. On the zero modes
 it acts as $p_L\mapsto p_L$, $p_R\mapsto-p_R$; on the field as $X=X_L+X_R\mapsto Y=X_L-X_R$, that is
 $\partial_\alpha X=\epsilon_{\alpha\beta}\partial^\beta Y$; on the oscillators as
 $\alpha_k\mapsto-\alpha_k$, $\tilde\alpha_k\mapsto\tilde\alpha_k$.
\item ``Symmetry'' has four levels: spectrum, conformal field theory, perturbation theory, exact
 theory. The mass formula shows the first; the invariance of all OPEs under a one-sided
 reflection shows the second; the path integral and the genus expansion show the third,
 provided $g_s\mapsto g_s\sqrt{\ap}/R$; the interpretation as a gauge transformation supports
 the fourth.
\item T-duality is a relabelling of one theory, not an invariance of each state, and not the
 statement that physics is independent of $R$.
\item The self-dual radius $R=\sqrt{\ap}$ is the unique fixed point. There, four extra massless
 vectors and eight extra massless scalars appear, $U(1)^2$ is enhanced to $SU(2)\times SU(2)$, and
 T-duality is the Weyl reflection of one $SU(2)$. The moduli space is the half-line
 $R\ge\sqrt{\ap}$.
\item Closed strings cannot measure a radius below $\sqrt{\ap}$: $R_{\rm obs}=\max(R,\ap/R)$. The
 smooth invariant $R_{\rm eff}=R+\ap/R\ge2\sqrt{\ap}$ is the basic object of the dual-scale
 programme; the inequality is \TA, its physical role is \TL{} or \TC{} as specified in
 \cref{sec:tduality:minlength}.
\item Lean has checked the zero-mode algebra over $\Q$ at $\ap=1$, the uniqueness of the fixed
 point, and $R+1/R\ge2$ over $\R$. It has not checked anything about oscillators, operator
 algebras, loops or the dilaton.
\end{itemize}
\end{summarybox}

\section*{Exercises}
\addcontentsline{toc}{section}{Exercises}

\begin{exercise}[$\star$]
Extend \cref{tab:tduality:example} to the circle of radius $R=3\ell_s$: list all label sets with
$\ap M^2\le1$, and write down the corresponding list at $\tilde R=\ell_s/3$ without further
computation. Which states are massless, and is either radius a point of enhanced gauge symmetry?
\end{exercise}

\begin{exercise}[$\star$]
Show that the map $(R,g_s)\mapsto(\ap/R,\;g_s\sqrt{\ap}/R)$ is an involution and that
$g_s^2/R$ and $\Phi-\frac12\ln(R/\sqrt{\ap})$ are invariant under it. For $R=\ell_s/5$ and
$g_s=0.01$, find the dual radius and coupling, and compare the two values of the combination
\eqref{eq:tduality:G25}.
\end{exercise}

\begin{exercise}[$\star$]\label{ex:tduality:special}
Besides the self-dual radius there are radii at which only extra massless \emph{scalars}
appear. With $N=\tilde N=0$, show that a pure winding state is massless when $R^2=4\ap/w^2$ and
a pure momentum state when $R^2=n^2\ap/4$ \source{Tong §8.2.3, ll.~11565--11590}. Check that these
two families of special radii are exchanged by T-duality, and explain why no massless charged
\emph{vector} can occur at any $R\ne\sqrt{\ap}$. (Hint: for $N=1$, $\tilde N=0$ show that
$M^2=p_R^2$.)
\end{exercise}

\begin{exercise}[$\star\star$]
Starting from $Y=X_L-X_R$, derive \eqref{eq:tduality:eps}. Show that under $X\to Y$ the
equation of motion and the identity $\partial_\tau\partial_\sigma X=\partial_\sigma\partial_\tau X$
are exchanged. Compute $\oint\dd\sigma\,\partial_\sigma Y$ and
$\frac1{2\pi\ap}\oint\dd\sigma\,\partial_\tau Y$ from \eqref{eq:tduality:Ymodes} and interpret the
results as winding and momentum on the dual circle.
\end{exercise}

\begin{exercise}[$\star\star$]
Prove the identity $\frac{\ap}2(p_Lp_L'-p_Rp_R')=nw'+wn'$ of \eqref{eq:tduality:VV}. Deduce
$h-\bar h=nw$ for $V_{n,w}$, and show that the bilinear form $\langle(n,w),(n',w')\rangle=nw'+wn'$
on $\Z^2$ is even, has signature $(1,1)$, and is preserved by exactly four integer matrices:
$\pm1$ and $\pm$ the swap. Which of these is T-duality, and what do the others do to $X_L$ and
$X_R$?
\end{exercise}

\begin{exercise}[$\star\star$]
Show that $R_{\rm eff}=2\sqrt{\ap}\cosh x$ and $\Tr\HH=2\cosh2x$, with
$x=\ln(R/\sqrt{\ap})$. Expand both to second order around the self-dual point. Show that any
smooth duality-invariant function $f(R)$ has $\dd f/\dd x=0$ at $x=0$, so that the self-dual
radius is a critical point of every duality-invariant potential for the radion. Why does this
not by itself show that the radius is stabilized there?
\end{exercise}

\begin{exercise}[$\star\star$, Lean]
In the namespace \lean{StringTheory.UseCases.TDuality}, state and prove that T-duality
exchanges the two $U(1)$ charges:
\begin{leancode}
theorem tduality_swaps_charges (n w R : ℚ) (hR : R ≠ 0) :
    leftMomentum w n (1 / R) + rightMomentum w n (1 / R)
      = leftMomentum n w R - rightMomentum n w R := by
  sorry
\end{leancode}
Give a proof that uses the two existing theorems and \lean{ring}, and a second one from the
definitions with \lean{field\_simp}. Which is more robust against a change of the definitions?
\end{exercise}

\begin{exercise}[$\star\star\star$, Lean]\label{ex:tduality:leansqrt}
State and prove the self-dual radius with $\ap$ restored:
for real $R>0$ and $\ap>0$ with $\ap/R=R$, one has $R=\sqrt{\ap}$ (\lean{Real.sqrt}). A possible
route: derive $R^2=\ap$ with \lean{field\_simp}, rewrite the goal with it, and use
\lean{Real.sqrt\_sq}. Is the hypothesis $\ap>0$ needed? Then explain why the analogous statement over $\Q$ cannot even be
formulated with a conclusion of the form $R=c$ for $\ap=2$.
\end{exercise}

\begin{exercise}[$\star\star\star$]
(Genus one by hand.) Let $\tau=\tau_1+\ii\tau_2$, $r=R/\sqrt{\ap}$, $q=\ee^{2\pi\ii\tau}$. Use the Poisson
formula $\sum_{m\in\Z}\ee^{-\pi a(m+b)^2}=a^{-1/2}\sum_{k\in\Z}\ee^{-\pi k^2/a+2\pi\ii kb}$ to show that
\[
  \sum_{m,n\in\Z}\exp\!\Big(-\frac{\pi r^2}{\tau_2}\,|m+n\tau|^2\Big)
  =\frac{\sqrt{\tau_2}}{r}\sum_{k,n\in\Z}q^{\frac12P_L^2}\,\bar q^{\frac12P_R^2},\qquad
  P_{L,R}=\frac1{\sqrt2}\Big(\frac kr\pm nr\Big),
\]
where $P_{L,R}=\sqrt{\ap/2}\,p_{L,R}$ are the dimensionless momenta of GPR. Including the factor
$r$ from the zero mode, conclude that $Z_1(1/r)=Z_1(r)$, the case $g=1$ of \eqref{eq:tduality:Zg}.
(The identity can be tested numerically, for instance at $r=1.3$, $\tau=0.23+0.9\,\ii$, where both
sides equal $1.01830\ldots$)
\end{exercise}

\section*{Further reading}
\addcontentsline{toc}{section}{Further reading}
\begin{itemize}[leftmargin=*]
\item The spectrum on a circle, enhanced symmetry, T-duality, the dual coordinate, the dilaton
 shift and the path-integral derivation: \source{Tong §8.2--8.3.1, ll.~11351--11803}.
\item Circle duality in the operator formalism, the one-loop and higher-genus partition
 functions and the dilaton: \source{GPR §2.2--2.3, ll.~576--1029}. Duality as a gauge symmetry
 and the Weyl reflection: \source{GPR §2.6.1, ll.~2377--2490}.
\item Abelian duality by gauging an isometry, global issues and the dilaton determinant:
 \source{AAL §2.1, ll.~215--376}; the operational definition of distance:
 \source{AAL §6, ll.~1289--1320}.
\item In this book: the general Buscher rules in \cref{ch:buscher}; the lattice of $(p_L,p_R)$ in
 \cref{ch:narain}; the group generated by dualities, shifts and basis changes in \cref{ch:odd};
 open strings and D-branes under T-duality in \cref{ch:dbranes}; the dual-scale bound for tori in
 \cref{ch:dualscale}.
\end{itemize}
